\section{From the arguments to Lean derivations}
\label{sec:lean}
Two theorem types play different roles here. A theorem about \texttt{Nat} in Lean checks a mathematical implementation in a familiar setting. A term of type $\Gamma\vdash[\mathsf{PRAAx}]A$ constructs a proof in the first-order object calculus. The input functions in a uniform object theorem are fixed PR definitions, or terms supplied by a metatheoretic schema; a general evaluator for all PR programs is unnecessary.

\subsection{The checked arithmetic components}
The archive has standalone Lean sources for trial division, rank selection, midpoint inversion, the factorial kernel and polynomial coefficients. They use Lean 4.19.0 and its standard library. In \texttt{Counting.lean}, \texttt{prime\_spec} identifies the Boolean test with the proper-divisor predicate. For any marked sequence, \texttt{select\_spec} proves that whenever the positive target rank has been reached by the supplied bound, the selected location is marked, has that exact rank, and has predecessor rank one less. The induction has two cases: the old prefix already reaches the target, or the new bit must be true and raise the count by precisely one. \texttt{select\_unique} establishes uniqueness of that marked location. The object proof still needs the Euclidean bound in Section~\ref{sec:found} to turn this bounded specification into unrestricted inverse rank.

\texttt{Inversion.lean} implements divisor-incidence inversion for integer sums. It proves the divisor sum of the recursively defined Möbius function, and proves the scaled \texttt{midpoint\_inversion} identity, parametrized by a Boolean marking function. Section~\ref{sec:powers} explains the correspondence of the scaled rows to $J_{1/2}$. The marking may be instantiated with \texttt{PrimeCounting.prime}. \texttt{ChebyshevKernel.lean} proves periodicity, the bit range, and the interval where the five-floor kernel equals one. \texttt{Growth.lean} verifies the two finite polynomial coefficient lists that decide the eventual sign of Ramanujan's inequality. \texttt{Sieve.lean} proves inclusion--exclusion for arbitrary lists of Boolean marks. Its \texttt{floor\_sieve} gives the alternating floor sum with subset lcm for arbitrary lists of moduli, including repetitions. For distinct primes the lcm equals the product by factorization. These files establish their indicated statements directly. They do not mechanize the analytic limits. See \texttt{verification/lean.txt} for the build and axiom output.

\subsection{The pinned object language}
Use ProveIt revision \texttt{a9c0f1949be61679ed2e219b97dc42ecb7ae91a9}, whose root project specifies Lean 4.32.0. The optional \texttt{ProveItPNT.lean} is a source-level client of the actual object theorem in \texttt{FinalPNT.lean}. Its type is the sequent of Section~\ref{sec:pntgrowth}. The included standalone verification used 4.19.0, so this client and the 4.32.0 ProveIt dependency closure have not been rebuilt here. The stronger logarithmic-integral error estimate and its $\vartheta$ version are reused from the separate human proof packet. Its authors describe its Lean material as an outline. Conclusions depending on these quantitative estimates carry that dependency even when their finite coefficient calculations have been checked.

In the following construction, \texttt{PTm n} is an arithmetic term in a context of $n$ number variables, and \texttt{Fm PTm n} is a formula. Compile every bounded scan into a primitive recursion storing its loop state and fixed parameters; project the needed component at the end. To use induction on a relation with a positive rational denominator, multiply both sides by the recorded denominators. Represent signed quantities by two nonnegative parts, and express $a-b\le c-d$ as $a+d\le c+b$. A finite interval or product can have its own denominator at each step; cross multiplication preserves order. The prime-power identity can instead use the common denominator $2(L!)^2$. Boolean universal and existence tests are justified by the prefix-product and first-hit invariants of Section~\ref{sec:found}.

The entries below give the proof obligations in dependency order. Module names in the last column are existing providers or statements to implement as indicated.
\begin{enumerate}
\item \textbf{Prime data.} Use \texttt{Valuation.isP} and the shipped theorems \texttt{isP\_of\_prime}, \texttt{prime\_of\_isP}, \texttt{isP\_01}. Use \texttt{PiTheta.Pi x zero} for the inclusive sum. Compile $E(k)$ and the bounded least scan. Express its minimality by the product bit for earlier failed comparisons. Obtain a successful prime search from the factorial argument before applying the rank lemma. Use \texttt{FactorTrace.factor\_product} or \texttt{FactorVector.reconstruct} for the positive factorization, and \texttt{factor\_injective} for reconstruction from equal valuation vectors.
\item \textbf{Survivor counts.} Define $\Phi$ directly by a double bounded sum/product. Define the subset mask product $d_s$ and parity fold. Prove the mask expansion by prefix induction in the integer-pair ring. Induct on $\lfloor N/d\rfloor$ and the final remainder to count the multiples. For $P_k$, compile the least-factor trace, its non-unit length and minimum index. Implement the two maps in \eqref{eq:sv-factor} by quotient and multiplication. To count carried indices, first sum their graph bits by source, then by target (finite double-sum swap). This yields the two explicit counts and both large sieve formulas. The wheel bound and first-crossing scan then need only signed inequalities. In the factorial congruence use the bounded inverse search modulo $p$ and pair off inverse residues.
\item \textbf{Weights and inversion.} Define the cut bits by $Bp^j<A$ and $Bp^j=A$; use the same $A,B$ throughout the $j=dk$ substitutions. Fill the Möbius table by increasing $n$, or use the positive and negative parts in the repository. Induct on $n$ to identify the divisor-recurrence version with the squarefree version, using their common divisor sum and isolating $d=n$. Existing helpers include \texttt{DivisorSwap.div\_swap}, \texttt{MobiusInversion.mobius\_inversion}, \texttt{PrimePowerMangoldt.lam\_pw} and \texttt{ChebyshevAbel.abel\_gen}. Transport the finite scaled inversion by clearing its denominator and swapping the two bounded sums. Form the lcm and prime product as integer codes before applying a log name.
\item \textbf{Analytic names.} Implement signed rational intervals for the four arithmetic operations. Inversion requires a named rational lower bound on the absolute denominator. Define $\log$ by rectangles or the range-reduced atanh series, $\exp$ by its factorial series, real powers by positive roots or $\exp(s\log x)$, and complex numbers by pairs of such enclosures. Define $\operatorname{li}$ by the finite exponential-integral expression in Section~\ref{sec:li-cut} with its tail. Define $\zeta$, $R$ and the step integrals by the cutoffs in Sections~\ref{sec:mellin}, \ref{sec:trivial} and \ref{sec:gram}. For each displayed equality, compile the finite expressions at requested precision $j+g$ on both sides and prove the width/error inequality for precision $j$. At a bounded fold fix the final rounding budget first, split it over all iterations, and induct on a numerical error accumulator. The circuit sensitivity bound and the explicit tail bounds supply $g$. This avoids induction on equality of real names, which would quantify over all precisions.
\item \textbf{Bounds and asymptotics.} Encode $\log t$, $\vartheta(t)$ and $\psi(t)$ by intervals with common requested accuracy. For the sharp global upper bound, prove the integer periodic kernel first, then the logarithm-of-factorial rectangle inequalities. Algebraically telescope the bounds for $\vartheta$ and use the finite data at $10^6$. For the PNT instantiate \texttt{arith\_pnt} and prove the change from harmonic mesh to the log name by the one-cell error bound. For higher estimates, the remaining substantial object-development dependency is translation of the stronger manuscript's finite Euler-product, contour and desmoothing estimates. Given those, the factorial polynomials, residual coefficient norms, eight finite improvements for $M$, and inverse substitutions in this report use the stated, bounded schedules. For $x/e$ use the factorial comparison in Section~\ref{sec:growthrefine}.
\item \textbf{Certificate reflection.} Compile the finite range checks as one state recursion: at each input update the prime count, lower/upper sums and a bit for all earlier comparisons. Prove once that its update preserves the intended arithmetic specification. For a fixed final height a numeral derivation reducing the bit to one furnishes a closed certificate. Large inputs are better replayed in blocks: store boundary states, verify each block and compose the state equalities by transitivity. The supplied Python/C++ runs evaluate the arithmetic; generating these object-calculus numeral derivations is further work.
\end{enumerate}

\paragraph{Proof scope.}
The full Riemann explicit formulas over the nontrivial zeros, their recommended summation orders and any interchange with the Möbius sum need analytic information beyond the Mellin identity in its half-plane of absolute convergence. This report proves the trivial-zero correction separately. Littlewood-type sign changes, zero-free regions at the sharper exponents quoted on Wikipedia, record-size table entries above $10^{14}$, and the published time and storage exponents of optimized sieves are outside the proved results. Section~\ref{sec:conditional-gap} specifies exactly which error bounds imply its short-interval conclusions. For non-global numerical lower bounds and estimates on $p_m$, the finite windows and the eventual results leave an intervening range unless a separate argument covers it.

\section*{Sources and reproduction}
\addcontentsline{toc}{section}{Sources and reproduction}
\begin{itemize}
\item \href{https://en.wikipedia.org/wiki/Prime-counting_function}{Wikipedia, \emph{Prime-counting function}}, and \href{https://mathworld.wolfram.com/PrimeCountingFunction.html}{Eric W. Weisstein, \emph{Prime Counting Function}, MathWorld}. The two requested entries identify the target formulas, endpoint conventions, numerical claims and tables, accessed November 22, 2026.
\item \href{https://github.com/VladimirReshetnikov/ProveIt/tree/a9c0f1949be61679ed2e219b97dc42ecb7ae91a9/Logic/PrimitiveRecursiveArithmetic}{Vladimir Reshetnikov, \emph{ProveIt}, pinned iPRA tree}. Relevant providers are the ConciseIPRA proof calculus, the ordinary derivations in \texttt{AnalyticNumberTheory}, and the \href{https://github.com/VladimirReshetnikov/ProveIt/tree/a9c0f1949be61679ed2e219b97dc42ecb7ae91a9/Logic/PrimitiveRecursiveArithmetic/docs/proofs/analytic-number-theory/pnt_li_ipra_2026-11-15}{logarithmic-integral packet}. The text identifies the two imported quantitative inputs separately.
\item \href{https://doi.org/10.1215/ijm/1255631807}{J. B. Rosser and L. Schoenfeld, \emph{Approximate formulas for some functions of prime numbers}}, Illinois J. Math. 6 (1962), 64--94. \href{https://doi.org/10.1007/s11139-016-9839-4}{P. Dusart, \emph{Explicit estimates of some functions over primes}}, Ramanujan J. 45 (2018), 227--251. Section~\ref{sec:chebsharp} gives an elementary constructive proof of the specific global upper bound.
\item \href{https://www-users.cse.umn.edu/~odlyzko/doc/arch/meissel.lehmer.pdf}{J. C. Lagarias, V. S. Miller and A. M. Odlyzko, \emph{Computing $\pi(x)$: the Meissel--Lehmer method}}, Math. Comp. 44 (1985), 537--560. \href{https://www.jstor.org/stable/2312732}{S. W. Golomb, \emph{On the ratio of $N$ to $\pi(N)$}}, Amer. Math. Monthly 69 (1962), 36--37.
\item \href{https://arxiv.org/abs/1407.1901}{A. Dudek and D. Platt, \emph{On solving a curious inequality of Ramanujan}}, Exp. Math. 24 (2015), 289--294. \href{https://arxiv.org/abs/1402.6417}{A. Dudek, \emph{On the Riemann hypothesis and the difference between primes}}. The conditional local weighted estimate is identified at \eqref{eq:local-triangle}.
\end{itemize}
